% texte.tex — f163r, vv. 43–79 (conversion initiale des travaux antérieurs)
\vers[n=43, indent=3em, findent=1.7999999999999998em]{}{}{ἐχέτω δὲ πνοὰς νήνεμος αἰθήρ,}{Que retienne ses haleines l'éther paisible,}
\vers[n=44, indent=3em, findent=1.7999999999999998em]{}{}{κῦμα δὲ πόντου μὴ κελαδείτω}{que le flot de la mer point ne retentisse}
\vers[n=45, indent=3em, findent=1.7999999999999998em]{}{}{γλαυκόν—}{glauque...}
\vers[suite]{ΜΝ.}{Mnésiloque}{βομβάξ.}{Bombax !}
\vers[suite]{ΕΥ.}{Euripide}{σίγα. τί λέγει;}{Chut ! que dit-il ?}
\vers[n=46, indent=3em, findent=1.7999999999999998em]{ΘΕ.}{Serviteur}{πτηνῶν τε γένη κατακοιμάσθω,}{Que la gent ailée s'assoupisse,}
\vers[n=47, indent=3em, findent=1.7999999999999998em]{}{}{θηρῶν τ᾽ ἀγρίων πόδες ὑλοδρόμων}{que les pieds des bêtes sauvages vaguant par les bois}
\vers[n=48, indent=3em, findent=1.7999999999999998em]{}{}{μὴ λυέσθων.}{point ne se meuvent...}
\vers[suite]{ΜΝ.}{Mnésiloque}{βομβαλοβομβάξ.}{Bombalo bombax !}
\vers[n=49, indent=3em, findent=1.7999999999999998em]{ΘΕ.}{Serviteur}{μέλλει γὰρ ὁ καλλιεπὴς Ἀγάθων}{... Car il s'apprête l'homme au beau langage, Agathon,}
\vers[n=50, indent=3em, findent=1.7999999999999998em]{}{}{πρανοῦν προτέρους—}{à aller au devant du gaillard...}
\vers[suite]{ΜΝ.}{Mnésiloque}{μῶν βινεῖσθαι;}{Et se faire enfiler ?}
\vers[n=51, indent=3em, findent=1.7999999999999998em]{ΘΕ.}{Serviteur}{τίς ὁ φωνήσας;}{Qui a parlé?}
\vers[suite]{ΜΝ.}{Mnésiloque}{νήνεμος αἰθήρ.}{L'éther paisible.}
\vers[n=52, indent=3em, findent=1.7999999999999998em]{ΘΕ.}{Serviteur}{δρυόχους, τιθέναι δράματος ἀρχάς.}{... du gaillard d'avant ; il va composer les débuts d'un drame.}
\vers[n=53, indent=3em, findent=1.7999999999999998em]{}{}{κάμπτει δὲ νέας ἁψῖδας ἐπῶν,}{Il cambre, en jantes nouvelles, des vers.}
\vers[n=54, indent=3em, findent=1.7999999999999998em]{}{}{τὰ δὲ τορνεύει, τὰ δὲ κολλομελεῖ,}{Il met les uns sur le tour, il assemble les autres.}
\vers[n=55, indent=3em, findent=1.7999999999999998em]{}{}{καὶ γνωμοτυπεῖ κἀντονομάζει}{Il frappe des sentences, fait s'opposer des mots,}
\vers[n=56, indent=3em, findent=1.7999999999999998em]{}{}{καὶ κηροχυτεῖ καὶ γογγύλλει}{malaxe la cire, l'entoure d'un moule,}
\vers[n=57, indent=3em, findent=1.7999999999999998em]{}{}{καὶ χοανεύει—}{puis introduit le godet...}
\vers[suite]{ΜΝ.}{Mnésiloque}{καὶ λαικάζει.}{Et avale.}
\vers[n=58, indent=3em, findent=1.7999999999999998em]{ΘΕ.}{Serviteur}{τίς ἀγροιώτας πελάθει θριγκοῖς;}{Quel est le manant qui approche cette enceinte ?}
\vers[n=59, indent=3em, findent=1.7999999999999998em]{ΜΝ.}{Mnésiloque}{ὃς ἕτοιμος σοῦ τοῦ τε ποιητοῦ}{Quelqu'un qui s'apprête à te mettre en boule et te retourner, toi et ton poète au beau langage,}
\vers[n=60, indent=3em, findent=1.7999999999999998em]{}{}{τοῦ καλλιεποῦς κατὰ τοῦ θριγκοῦ}{et à introduire par votre enceinte}
\vers[n=61, indent=3em, findent=1.7999999999999998em]{}{}{συγγογγύλας καὶ συστρέψας}{}
\vers[n=62, indent=3em, findent=1.7999999999999998em]{}{}{τουτὶ τὸ πέος χοανεῦσαι.}{cet outil-ci.}
\vers[n=63]{ΘΕ.}{Serviteur}{ἦ που νέος γ᾽ ὢν ἦσθ᾽ ὑβριστὴς ὦ γέρον.}{Dans ta jeunesse tu devais être un ribaud, vieillard..}
\vers[n=64]{ΕΥ.}{Euripide}{ὦ δαιμόνιε τοῦτον μὲν ἔα χαίρειν, σὺ δὲ}{(À Mnésiloque :) Diantre d'homme, laisse-le tranquille. (Au serviteur :) Toi,}
\vers[n=65]{}{}{Ἀγάθωνά μοι δεῦρ᾽ ἐκκάλεσον πάσῃ τέχνῃ.}{va m'appeler Agathon bon gré, mal gré.}
\vers[n=66]{ΘΕ.}{Serviteur}{μηδὲν ἱκέτευ᾽· αὐτὸς γὰρ ἔξεισιν τάχα.}{Inutile d'insister. Lui-même ne tardera pas à sortir,}
\vers[n=67]{}{}{καὶ γὰρ μελοποιεῖν ἄρχεται· χειμῶνος οὖν}{car il commence un chant lyrique. Et de fait, en hiver,}
\vers[n=68]{}{}{ὄντος κατακάμπτειν τὰς στροφὰς οὐ ῥᾴδιον,}{il n'est pas facile de contourner ses strophes,}
\vers[n=69]{}{}{ἢν μὴ προίῃ θύρασι πρὸς τὸν ἥλιον.}{sans aller se mettre dehors au soleil.}
\vers[n=70]{ΜΝ.}{Mnésiloque}{τί οὖν ἐγὼ δρῶ;}{Et moi, que vais-je faire ?}
\vers[suite]{ΕΥ.}{Euripide}{περίμεν᾽, ὡς ἐξερχεται.}{Demeure, puisqu'il arrive.}
\vers[n=71]{ΜΝ.}{Mnésiloque}{ὦ Ζεῦ τί δρᾶσαι διανοεῖ με τήμερον;}{Ô Zeus, que prétends-tu faire de moi aujourd'hui ?}
\vers[n=72]{}{}{νὴ τοὺς θεοὺς ἐγὼ πυθέσθαι βούλομαι}{Par les dieux, je tiens à savoir de quoi il retourne.}
\vers[n=73]{}{}{τί τὸ πρᾶγμα τουτί. τί στένεις; τί δυσφορεῖς;}{Qu'as-tu à geindre ? à te contrarier ?}
\vers[n=74]{}{}{οὐ χρῆν σε κρύπτειν ὄντα κηδεστὴν ἐμόν.}{Il n'y a pas de mystère à faire avec moi, puisque tu es mon gendre.}
\vers[n=75]{ΕΥ.}{Euripide}{ἔστιν κακόν μοι μέγα τι προπεφυραμένον.}{Un grand malheur, manigancé d'avance, me menace.}
\vers[n=76]{ΜΝ.}{Mnésiloque}{ποῖόν τι;}{De quelle sorte ?}
\vers[suite]{ΕΥ.}{Euripide}{τῇδε θἠμέρᾳ κριθήσεται}{Ce jour décidera}
\vers[n=77]{}{}{εἴτ᾽ ἔστ᾽ ἔτι ζῶν εἴτ᾽ ἀπόλωλ᾽ Εὐριπίδης.}{si Euripide vit encore ou s'il est mort.}
\vers[n=78]{ΜΝ.}{Mnésiloque}{καὶ πῶς; ἐπεὶ νῦν γ᾽ οὔτε τὰ δικαστήρια}{Comment cela ? Car ni les tribunaux}
\vers[n=79]{}{}{μέλλει δικάζειν οὔτε βουλῆς ἐσθ᾽ ἕδρα,}{ne sont dans le cas de juger, ni la Boulé ne siège,}
